% TOP_PROBLEM: {"id": 3384, "title": "Davenport's constant", "category_id": 1, "category": {"id": 1, "name": "number_theory", "display_name": "Number Theory", "description": "Properties of integers, prime numbers, Diophantine equations.", "slug": "number-theory", "order_index": 1, "created_at": "2026-07-31T15:26:25.670Z"}, "status": "open", "problem_number": "OPG-563", "set_id": 12}
% FIRST_POSTED: 2026-10-01
% ENTRY_KIND: statement_only
% UnsolvedMath ID 3384; source catalogue code OPG-563
% !TeX program = lualatex
% Definition and sources only; this file is not an LLM solution attempt.
\documentclass[11pt,a4paper]{article}
\usepackage[margin=25mm]{geometry}
\usepackage{fontspec}
\setmainfont{lmroman10-regular.otf}[BoldFont=lmroman10-bold.otf,
 ItalicFont=lmroman10-italic.otf,BoldItalicFont=lmroman10-bolditalic.otf]
\usepackage{amsmath,amssymb,mathtools}
\usepackage{unicode-math}
\setmathfont{latinmodern-math.otf}
\AtBeginDocument{\let\setminus\smallsetminus}
\usepackage{xurl,hyperref}
\hypersetup{colorlinks=true,urlcolor=blue}
\setcounter{secnumdepth}{0}
\setlength{\parindent}{0pt}
\setlength{\parskip}{5pt}
\setlength{\emergencystretch}{3em}
\begin{document}
\section*{Davenport's constant}\label{3384}
{\small UnsolvedMath ID 3384 · OPG-563 · Number Theory · OpenGarden}
\subsection{Definitions and mathematical statement}
For a finite abelian group $G$ written additively, the Davenport
constant $D(G)$ is the least positive integer $m$ such that every
sequence $(g_1,\ldots,g_m)$ in $G$ has a nonempty set of positions
$I$ satisfying $\sum_{i\in I}g_i=0$. Terms may repeat, and positions
may be selected without being consecutive. This constant is denoted
$s(G)$ in the source; the symbol $D$ avoids confusion with a
fixed-length zero-sum threshold.

For integers $n\ge2$ and $d\ge1$, the conjectural formula is
\[
              D((\mathbb Z/n\mathbb Z)^d)=d(n-1)+1.
\]
Here $d$ is the number of cyclic direct factors, all with the same
modulus $n$. The lower bound is witnessed by the sequence containing
$n-1$ copies of each standard coordinate vector: a nonempty selection
cannot sum to zero because each coordinate count lies between zero
and $n-1$. Consequently the target is the upper bound asserting a
zero sum in every sequence of length $d(n-1)+1$.
The quantifier covers every modulus and every dimension, without
restricting $n$ to a prime or prime power. A statement for subsets
with no repeated elements would be a different extremal problem.
\subsection{Short English statement}
Must every sequence of d(n minus one) plus one elements in d copies of the cyclic group of order n contain a nonempty zero-sum selection?
\subsection{Sources}
\begin{itemize}
\item[S1] ULAM.AI and the UnsolvedMath Contributors, supplied problems.json, record 3384 (OPG-563), OpenGarden collection. Catalogue identity and source transcription; CC BY 4.0.
\url{https://huggingface.co/datasets/ulamai/UnsolvedMath}.
\item[S2] Open Problem Garden, Davenport's constant. Original problem and discussion; the mathematical formulation below is expanded from the supplied source record.
\url{http://www.openproblemgarden.org/op/davenports_constant}.
\end{itemize}
\end{document}
